{\small\textbf{The one buffer built to be released (Zheng, 1:00) [2:30–3:30]}}\par\smallskip
Why does it matter \emph{which} buffer you hold? Because not all capital is usable in a crisis.
\begin{itemize}\setlength\itemsep{0pt}
\item \textcolor{navy}{\textbf{[def]}} The countercyclical capital buffer, the CCyB, is the only layer of this stack that the authority can \textbf{switch off} overnight.
\item It ranges from 0 to 2.5\% of risk-weighted assets. \textcolor{navy}{\textbf{[def]}} Those are a bank's assets weighted by their riskiness, so a safe mortgage counts for less than a risky corporate loan.
\item Increases take a year to bind; releases are immediate.
\end{itemize}
\par\smallskip
The other buffers are technically usable, but if a bank dips into them, it faces automatic limits on dividends and bonuses. \textcolor{navy}{\textbf{[def]}} That is the "maximum distributable amount", or MDA, restriction. Banks hate that, and markets punish it.
\begin{itemize}\setlength\itemsep{0pt}
\item ECB research shows that in COVID, banks close to these thresholds \textbf{cut lending}, even though they were allowed to use the buffers.
\item Only a buffer the authority explicitly releases is truly usable.
\end{itemize}
\par\smallskip
That is the lesson of COVID. What matters is not only how much capital banks hold, but how much of it the authority can actually hand back.
